% ============================================================================
% Fichier  : partie1/P01_C02_droit.tex
% Titre    : Le Cadre Juridique Camerounais
% Auteur   : MINKA MI NGUIDJOÏ Thierry Emmanuel
% Version  : 1.0 -- novembre 2026
% ============================================================================
% Sources : loi n° 2010/012 du 21 décembre 2010 (texte officiel : art. 65,
% 66, 68, 69, 70, 71) ; loi n° 2024/017 (principes seulement).
% ============================================================================

\chapter{Le Cadre Juridique Camerounais}
\label{chap:droit}
\chaptermark{Le cadre juridique}

\epigraph{%
    « La curiosité n'est pas une autorisation. »%
}{Principe du cours}

\niveauChapitre{Fondamental}{Chapitre~\ref{chap:ethique}~; cours SEC3266,
chapitre sur le droit.}

\begin{boxobjectifs}
\begin{itemize}
    \item Citer les infractions de la loi n°~2010/012 qui concernent le
          testeur.
    \item Expliquer pourquoi l'autorisation écrite fait la différence entre
          un test et une infraction.
    \item Identifier qui peut autoriser un test, et qui ne le peut pas.
    \item Repérer les risques juridiques propres aux données personnelles,
          aux tiers et aux outils.
    \item Adopter les réflexes qui protègent le testeur et son client.
\end{itemize}
\end{boxobjectifs}

\bigskip

% ============================================================================
\section*{Question directrice : qu'est-ce qui sépare Mme NGO MBOCK de l'attaquant de mars~?}
% ============================================================================

En mars, un attaquant est entré dans le système de LiZbiZ, s'y est
maintenu et a obtenu les droits d'administrateur. En octobre, Mme~NGO~MBOCK
a fait presque la même chose. Les gestes se ressemblent~; les qualifications
juridiques sont opposées. La différence tient en une feuille signée. Ce
chapitre en explique la portée.

\medskip
\textbf{Plan du chapitre.}
\begin{enumerate}
    \item Les infractions qui concernent le testeur (section~\ref{sec:c02-infractions}).
    \item L'autorisation (section~\ref{sec:c02-autorisation}).
    \item Les données, les tiers, les outils (section~\ref{sec:c02-risques}).
    \item Les réflexes du testeur (section~\ref{sec:c02-reflexes}).
\end{enumerate}

\separateur

% ============================================================================
\section{Les Infractions qui Concernent le Testeur}
\label{sec:c02-infractions}
% ============================================================================

\begin{boxjuris}
\textbf{\loicam{n°~2010/012} du 21~décembre 2010 relative à la
cybersécurité et la cybercriminalité au Cameroun} (texte officiel)

\medskip
\textbf{Article~65, alinéa~2}~: \textit{« Est puni des peines prévues à
l'alinéa 1 ci-dessus, tout accès non autorisé, à l'ensemble ou à une partie
d'un réseau de communications électroniques ou d'un système d'information
ou d'un équipement terminal. »}

\medskip
\textbf{Article~65, alinéa~4}~: \textit{« Est puni des mêmes peines prévues
à l'alinéa 1 ci-dessus, celui qui, sans droit, permet l'accès dans un réseau
de communications électroniques ou dans un système d'information par défi
intellectuel. »}

\medskip
\textbf{Portée.} Les peines de l'alinéa~1 sont de cinq à dix ans
d'emprisonnement et de 5 à 10~millions de FCFA d'amende, ou l'une de ces
deux peines. L'alinéa~4 vise expressément le « défi intellectuel »~: la
bonne intention ou la curiosité ne justifient pas un accès sans droit.
\end{boxjuris}

\begin{tableaucours}{L{2.0cm}L{6.0cm}L{4.2cm}}{%
    \tableauHeader{Article} &
    \tableauHeader{Comportement visé} &
    \tableauHeader{Geste de test concerné}
}
65 (1) & Interception de données sans droit ni autorisation & Écoute du
réseau, interception d'échanges \\
65 (2) et (4) & Accès non autorisé~; accès permis par défi intellectuel &
Toute connexion à un système hors mandat \\
66 & Perturbation ou interruption du fonctionnement par l'altération des
données & Test qui fait tomber un service \\
68 & Accès ou maintien frauduleux avec perturbation grave & Exploitation
qui endommage \\
69 & Accès sans droit, en violation des mesures de sécurité, pour obtenir
des données & Lecture de données au-delà du nécessaire \\
70 & Attaque par saturation pour faire effondrer une ressource & Test de
charge ou de déni de service \\
71 & Introduction de données pour supprimer ou modifier les données &
Modification de données de production \\
\end{tableaucours}
\vspace{-0.8em}
\begin{center}{\small\itshape Tableau~2.1. Les articles de la loi n°~2010/012 qui concernent le testeur}\end{center}

% ============================================================================
\section{L'Autorisation}
\label{sec:c02-autorisation}
% ============================================================================

L'autorisation écrite rend l'accès « autorisé » au sens de l'article~65.
Encore faut-il qu'elle émane de la personne qui a le pouvoir de la
donner~: le propriétaire ou l'exploitant du système, par son représentant
habilité.

\begin{tableaucours}{L{4.4cm}L{7.8cm}}{%
    \tableauHeader{Qui peut autoriser~?} &
    \tableauHeader{Commentaire}
}
Le dirigeant de l'organisation & Oui, ou son délégataire désigné par écrit \\
Le responsable informatique & Seulement s'il a reçu délégation \\
Un salarié qui « veut savoir » & Non \\
L'hébergeur ou le fournisseur de nuage & Pour son infrastructure~; le
client peut avoir besoin de son accord selon le contrat \\
Un tiers dont les systèmes sont reliés & Non pour le système du client~;
il autorise seulement ses propres systèmes \\
\end{tableaucours}
\vspace{-0.8em}
\begin{center}{\small\itshape Tableau~2.2. Le pouvoir d'autoriser}\end{center}

\begin{boxalert}
\textbf{Violation typique.} Se contenter d'un accord oral ou d'un courriel
d'un technicien.

\smallskip
\textbf{Règle~:} l'autorisation est écrite, signée par une personne
habilitée, datée, et précise le périmètre et les méthodes. Le testeur en
garde une copie sur lui pendant les tests physiques et dans son dossier.
\end{boxalert}

% ============================================================================
\section{Les Données, les Tiers, les Outils}
\label{sec:c02-risques}
% ============================================================================

\begin{boxinfo}
\textbf{Les données personnelles.} Un test peut faire voir des données
personnelles (fichier des clients, paie). La loi n°~2024/017 relative à la
protection des données à caractère personnel s'applique à leur traitement~;
ses articles sont à lire dans le texte publié au Journal officiel. La règle
du cours~: ne consulter que ce qui prouve la faille, ne rien copier,
masquer les preuves, détruire ce qui reste à la fin de la mission.
\end{boxinfo}

\begin{tableaucours}{L{3.0cm}L{9.2cm}}{%
    \tableauHeader{Risque} &
    \tableauHeader{Parade}
}
Tiers touché & Vérifier la propriété de chaque adresse avant de tester~;
exclure ce qui appartient à d'autres \\
Interruption de service & Exclure le déni de service~; prévenir avant
toute action risquée~; mot d'arrêt \\
Outils & Utiliser des outils de provenance sûre, sur des postes dédiés~;
ne pas diffuser d'outil offensif hors cadre \\
Preuves & Tracer chaque action (horodatage, adresse source) pour pouvoir
prouver ce qui a été fait, et ce qui ne l'a pas été \\
\end{tableaucours}
\vspace{-0.8em}
\begin{center}{\small\itshape Tableau~2.3. Risques juridiques et parades}\end{center}

\begin{boxscenario}{Opération LOGONE. L'incident évité}
\textbf{Constatation.} Le 26~octobre, M.~TAMO s'apprête à balayer la plage
10.10.0.0/16 entière. Mme~NGO~MBOCK lui rappelle l'exclusion de
10.10.50.0/24 (paie, hébergée pour un tiers). La commande est corrigée avant
d'être lancée~; l'incident est consigné.

\medskip
\textbf{Analyse} (confiance élevée). Le balayage aurait touché un système
dont LiZbiZ n'est pas propriétaire~: un accès non autorisé au sens de
l'article~65, et une faute contractuelle du cabinet.
\end{boxscenario}

% ============================================================================
\section{Les Réflexes du Testeur}
\label{sec:c02-reflexes}
% ============================================================================

\begin{boxprocedure}
\textbf{Avant, pendant, après}
\begin{enumerate}
    \item \textbf{Avant}~: lettre signée par une personne habilitée~;
          périmètre vérifié adresse par adresse~; contrat avec clause de
          responsabilité et assurance du cabinet.
    \item \textbf{Pendant}~: rester dans le périmètre et la fenêtre~;
          journaliser chaque action~; s'arrêter et prévenir au moindre
          doute ou incident.
    \item \textbf{Après}~: remettre le rapport au seul client~; détruire les
          données et les accès créés~; conserver le journal des actions.
\end{enumerate}
\end{boxprocedure}

\begin{boxmetier}{Auditeur de sécurité}
\textbf{Vérifier le mandat avant le premier geste.} L'auditeur qui intègre
des tests techniques vérifie lui-même la lettre, le périmètre et le
pouvoir du signataire~; il ne se fie pas à ce que lui dit l'audité.
\end{boxmetier}

\separateur

\begin{boxresume}
\textbf{Ce qu'il faut retenir de ce chapitre~:}
\begin{itemize}
    \item \textbf{L'article~65} réprime l'accès non autorisé, y compris
          « par défi intellectuel ».
    \item \textbf{Les articles~66 à~71} visent la perturbation, le maintien
          frauduleux, l'obtention de données, la saturation, la modification.
    \item \textbf{L'autorisation} est écrite, signée par une personne
          habilitée, précise.
    \item \textbf{Données, tiers, outils, preuves}~: quatre risques à
          maîtriser.
    \item \textbf{Avant, pendant, après}~: les réflexes qui protègent le
          testeur et le client.
\end{itemize}
\end{boxresume}

\begin{boxinfo}
\textbf{Points de contrôle de l'élève.} Je sais~: citer l'article~65 et ses
alinéas~2 et~4~; relier un geste de test à un article~; dire qui peut
autoriser~; énumérer les réflexes avant, pendant et après.
\end{boxinfo}

\bigskip

\begin{boxexo}[title={\ding{45}\quad Exercice~2.1. Quel article~? \niveaubase}]
Quel article vise~: un test de charge qui fait tomber un site~; la lecture
de données clients au-delà du nécessaire~; un accès « pour voir »~?
\end{boxexo}

\begin{boxexo}[title={\ding{45}\quad Exercice~2.2. Défi intellectuel \niveaubase}]
Que signifie, pour un étudiant passionné, l'alinéa~4 de l'article~65~?
\end{boxexo}

\begin{boxexo}[title={\ding{45}\quad Exercice~2.3. Les peines \niveaubase}]
Quelles sont les peines prévues à l'article~65, alinéa~1~?
\end{boxexo}

\begin{boxexo}[title={\ding{45}\quad Exercice~2.4. Qui autorise~? \niveaumoyen}]
Un administrateur de la BANQUE SANAGA vous envoie un courriel « vous pouvez
tester notre serveur ». Que faites-vous~?
\end{boxexo}

\begin{boxexo}[title={\ding{45}\quad Exercice~2.5. Le nuage \niveaumoyen}]
LiZbiZ veut faire tester sa messagerie en ligne. Quelles autorisations
faut-il vérifier~?
\end{boxexo}

\begin{boxexo}[title={\ding{45}\quad Exercice~2.6. Les données \niveaumoyen}]
Pendant un test, vous accédez au fichier des clients. Que faites-vous,
étape par étape~?
\end{boxexo}

\begin{boxexo}[title={\ding{45}\quad Exercice~2.7. L'incident évité \niveaumoyen}]
Rédigez l'entrée du journal de mission sur l'incident évité du 26~octobre.
\end{boxexo}

\begin{boxexo}[title={\ding{45}\quad Exercice~2.8. L'outil \niveaumoyen}]
Un collègue veut publier sur Internet l'outil qu'il a écrit pendant une
mission. Quels risques juridiques et contractuels voyez-vous~?
\end{boxexo}

\begin{boxexo}[title={\ding{45}\quad Exercice~2.9. Le contrat \niveauexpert}]
Rédigez les clauses juridiques d'un contrat de test d'intrusion~:
autorisation, responsabilité, assurance, données, confidentialité.
\end{boxexo}

\begin{boxexo}[title={\ding{45}\quad Exercice~2.10. Comparer \niveauexpert}]
Comparez la situation de l'attaquant de mars et celle de Mme~NGO~MBOCK
au regard des articles~65, 68 et~69~: éléments matériels, autorisation,
conséquences.
\end{boxexo}

\bigskip

\begin{boxinfo}
\textbf{Pour aller plus loin.}
\begin{itemize}
    \item Loi n°~2010/012~\cite{loi2010012}~; loi n°~2024/017~\cite{loi2024017}.
\end{itemize}
\end{boxinfo}

\begin{boxinfo}
\textbf{Transition.} L'autorisation se traduit en règles concrètes~: les
règles d'engagement, objet du Chapitre~\ref{chap:engagement}.
\end{boxinfo}
